

\section{Future Work}

This study provides valuable insights into the application of the Contrast Avoidance Model (CAM) using social media data, yet several avenues for future research remain to be explored. First, expanding the sample size and diversity of the dataset could help improve the generalizability of the findings. Future studies should seek to incorporate data from additional social media platforms beyond Twitter (e.g., Facebook, Instagram, Reddit) and consider non-English language content to account for cultural and linguistic differences in affective expression. This would provide a more comprehensive understanding of how CAM-related affective patterns manifest across different digital contexts.

Another critical avenue for future work is the clinical validation of user cohorts. In this study, anxiety and depression were identified through self-reported diagnoses via Twitter, which may not be as reliable as clinically diagnosed conditions. Collaboration with mental health professionals or the use of clinical survey data would help confirm the diagnostic accuracy of cohort membership, thus enhancing the validity of the conclusions drawn.

Moreover, while VADER-based sentiment analysis is useful in capturing continuous sentiment, it is limited by its reliance on lexicons and predefined rules. While a transformer based approach is presented in this study, its findings opposed the previous findings dramatically. Since the sample size for the transformer-based sentiment scores is relatively low, future work should attempt to validate the findings of the present study. If the results can be confirmed, that would imply, that the CAM cannot be easily validated empirically, or even present arguments invalidating the CAM to a certain degree.

Furthermore, while the proposed LMM-based analysis provides a significantly cleaner perspective on NA variability differences by accounting for covariates, more external influential factors should be included. In future studies, an exact age of the user could be added. Although the age groups presented in this study are useful, they might lack the granularity in order to confidently predict NA levels in users. The time of disease onset would also be an important predictor, together with possible therapeutic interventions that the user has used. Getting a more comprehensive representation of the user's disease progression state is highly likely to be an influential factor in NA level prediction.

In addition, expanding the scope of sentiment analysis to include multimodal signals — such as images, emojis, hashtags, or user interactions — would offer a richer, more nuanced view of affective expression. Social media users frequently convey emotion not only through text but also through visual elements and engagement patterns. Incorporating these multimodal sources of data could improve the accuracy of affective measurements and enable more robust analyses of emotional dynamics. This would also include hyperlink, which are included in many of the analyzed tweets, but their referenced were not resolved in this study. They might contain useful information to more accurately estimate NA levels.

Temporal modeling is another important area for development. While this study employed windowed aggregation of sentiment data to capture temporal variability in NA, future work could explore more sophisticated methods for modeling affect over time. Techniques such as dynamic structural equation modeling (DSEM, \citet{dsem})  could offer new ways to examine how affective states evolve within individuals and across cohorts, potentially revealing more complex patterns of emotional regulation or instability.

Finally, extending the application of CAM to additional clinical or subclinical groups represents an exciting direction for future research. By broadening the scope to include individuals with conditions such as post-traumatic stress disorder (PTSD), bipolar disorder, or generalized emotional distress, future studies could further validate the universality of CAM in capturing emotional dynamics across diverse mental health conditions.

In sum, while this study provides important initial evidence supporting the CAM framework in social media research, there are numerous opportunities to refine the methodology, expand the dataset, and incorporate advanced analytical techniques. These efforts will help enhance the robustness of the findings and facilitate the broader application of CAM in understanding emotional processes in digital environments.